% ------------------------------------------------------------------
% Introduction (no subsections; published literature only)
% ------------------------------------------------------------------

Lower respiratory tract infection is frequently treated without a microbiological diagnosis. In a large prospective study of adults admitted with community-acquired pneumonia, no pathogen was identified in 62\% despite extensive conventional testing\cite{jain2015community}. In critically ill patients, previous antibiotic treatment further reduces culture sensitivity\cite{miao2018microbiological}, noninfectious inflammatory syndromes can resemble pneumonia\cite{langelier2018integrating}, and the differential diagnosis extends to fungi, viruses and opportunistic pathogens. Metagenomic next-generation sequencing (\mngs) was developed to address this gap by detecting bacterial, fungal, viral and parasitic nucleic acids without a prespecified target\cite{chiu2019clinical,simner2018understanding}. Clinical studies have shown its diagnostic potential in meningoencephalitis\cite{wilson2019clinical}, bloodstream infection\cite{blauwkamp2019analytical,hogan2021clinical} and lower respiratory tract infection\cite{charalampous2019nanopore,langelier2018integrating,chen2020clinical,diao2022metagenomics}, and the assay is increasingly used in critically ill patients in Asia\cite{miao2018microbiological,han2019mngs}. Its breadth, however, carries a clinical risk: tests that detect incidentally carried microbes can lead to antimicrobial overuse and adverse outcomes\cite{mick2023integrated}.

The main challenge is interpretation rather than detection. Respiratory \mngs identifies upper-airway flora\cite{charlson2011topographical,dickson2016microbiome}, reagent and environmental contaminants\cite{salter2014reagent}, herpesviruses shed or reactivated during critical illness\cite{luyt2007herpes,limaye2008cytomegalovirus} and airway-colonizing \textit{Candida}\cite{meersseman2009significance}, as well as causative pathogens. In a prospective bronchoalveolar-lavage cohort, \mngs identified a definite or probable pathogen in 65\% of patients with lower respiratory tract infection, compared with 20\% by culture, but had a specificity of 75.4\% and a positive predictive value of 78.5\% against a composite reference standard\cite{chen2020clinical}. Distinguishing infection from colonization therefore remains a central limitation\cite{simner2018understanding,diao2022metagenomics}. Read thresholds and subtraction of no-template controls reduce background signal but do not resolve the biological ambiguity of a colonized respiratory specimen\cite{miller2019laboratory,schlaberg2017validation}.

Resolving that ambiguity requires evidence from outside the sequencing assay. Culture establishes viability and enables susceptibility testing\cite{metlay2019diagnosis,kalil2016management}, multiplex PCR provides rapid, validated target-specific results\cite{buchan2020practical}, galactomannan contributes an independent criterion for invasive fungal disease\cite{donnelly2020revision}, serology can establish infection after the organism has cleared from the airway\todo{cite}, and imaging defines the anatomical syndrome of pneumonia\cite{metlay2019diagnosis}. Each result must be read against the specimen and its timing: infection confined to blood, tissue, an abscess or an empyema may be absent from respiratory fluid, and intermittent or treated infection may be missed outside its detection window\cite{gu2019clinical,miller2022role,schlaberg2017validation}. Determining causation consequently requires integration of sequencing, conventional microbiology and clinical context.

Computational approaches to this problem have so far drawn only on the sequencing data. A rules-based model that scored organisms by relative abundance and known pathogenicity separated pathogens from commensals in tracheal aspirates from critically ill adults\cite{langelier2018integrating}, and classifiers adding host gene expression improved diagnosis of lower respiratory tract infection in critically ill children\cite{mick2023integrated} and of sepsis from plasma\cite{kalantar2022integrated}. Their host-response components require RNA sequencing of host transcripts, and none incorporates culture, antigen or serological results or the clinical course that physicians use to adjudicate causation.

Large language models could perform this integration because they encode clinical knowledge\cite{singhal2023large}, generate differential diagnoses from vignettes\cite{mcduff2025towards,tu2025towards} and have been evaluated as infectious-diseases consultants\cite{maillard2024chatbot}. Their performance on curated tasks, however, has not consistently transferred to clinical records. In an evaluation of 2,400 patient cases, current models performed worse than physicians, did not reliably follow diagnostic or treatment guidelines, had difficulty interpreting laboratory results and were sensitive to both the quantity and the order of the information they received\cite{hager2024evaluation}. Model outputs also vary between runs and may not reveal which evidence determined a prediction\cite{ji2023survey}. A generated explanation cannot by itself establish fidelity because a model's stated reasoning may not correspond to the computation that produced its answer\cite{turpin2023language}; in clinical applications, such post-hoc explanations can provide false reassurance\cite{ghassemi2021false}. Multi-agent systems\cite{tang2024medagents,kim2024mdagents,qiu2024llm} and retrieval augmentation\cite{lewis2020retrieval,xiong2024benchmarking} address parts of this problem, but clinical use also requires reproducible decisions, traceable evidence and an explanation written separately from the prediction. Evaluation must also guard against leakage\cite{kapoor2023leakage}: clinical notes often record the treating team's final diagnosis, which a model can copy rather than infer. To our knowledge, no approach has combined \mngs with conventional microbiology and the clinical record to attribute causation organism by organism under these constraints.

Here we present \toolname (\toolexpansion), which integrates respiratory \mngs with conventional microbiology and clinical data. Modality-specific language-model agents extract structured evidence; deterministic rules select candidate pathogens; and case-matched PubMed evidence can rescue candidates not selected initially. The clinical rationale is written only after the selected list has been fixed. In patients with suspected lower respiratory tract infection from three medical centres in Taiwan, we compared \toolname with individual diagnostic read-outs and with direct prompting of general-purpose language models against physician-adjudicated pathogens. \toolname identified the adjudicated pathogens more accurately than any single read-out or directly prompted model, reported far fewer organisms than the sequencing report while retaining most true pathogens, and linked each selection and exclusion to the evidence behind it. The system is intended to provide an evidence-linked candidate list for physician review, not an autonomous diagnosis.
